\documentclass[a4paper]{article}
\usepackage[margin=1.5cm]{geometry}
\usepackage{graphicx}
\usepackage{subcaption}
\usepackage{caption}
\usepackage{booktabs}
\usepackage{array}

\captionsetup[figure]{skip=2pt}

\newcommand{\schemefig}[1]{%
  \begin{minipage}[t]{0.49\linewidth}
    \centering
    \includegraphics[width=\linewidth]{outputs/#1/velocity_profile.pdf}\\[2pt]
    {\small\bfseries #1}
  \end{minipage}}

\begin{document}

\begin{center}
  {\LARGE\bfseries Gresho vortex -- Velocity norm vs radius}
\end{center}

\bigskip

\noindent
\textbf{Test case.}
Gresho vortex on a $20\times20$ quad mesh, $[0,1]^2$ domain, center at $(0.5,0.5)$.
Mach number $M_\infty = 10^{-5}$ (very low Mach, compressible solver).
CFL$=0.4$, $t_\mathrm{max}=10^{-4}$.
Analytical solution: $v_\theta = 5r$ for $r < 0.2$, $v_\theta = 2 - 5r$ for $0.2 \le r < 0.4$, $v_\theta = 0$ for $r \ge 0.4$.

\medskip
\noindent
\textbf{Note.}
The large gap between numerical and analytical results is expected at first order on this
coarse mesh in the very-low-Mach regime: compressible solvers exhibit $O(M^{-2})$ excess
dissipation, which effectively damps the vortex over $t_\mathrm{max}=10^{-4}$.

\bigskip

%% --- Row 1: classical schemes ---
\noindent
\schemefig{two\_wave}\hfill\schemefig{three\_wave}

\noindent
\schemefig{modified\_three\_wave}\hfill\schemefig{multi\_point}

\noindent
\schemefig{multi\_point\_iso}\hfill\phantom{\schemefig{multi\_point\_iso}}

\newpage

%% --- Row 2: ZB AR1D ---
\noindent
\schemefig{ZB\_AR1D\_LS}\hfill\schemefig{ZB\_AR1D\_LSM}

\medskip\noindent
\schemefig{ZB\_ARMD\_LS}\hfill\schemefig{ZB\_ARMD\_LSM}

\medskip\noindent
\schemefig{ZB\_ARMDM\_LS}\hfill\schemefig{ZB\_ARMDM\_LSM}

\medskip\noindent
\schemefig{ZB\_ARMDMAT\_LS}\hfill\schemefig{ZB\_ARMDMAT\_LSM}

\newpage

%% --- Row 3: ZB ARMDMMAT + AM + AMISO ---
\noindent
\schemefig{ZB\_ARMDMMAT\_LS}\hfill\schemefig{ZB\_ARMDMMAT\_LSM}

\medskip\noindent
\schemefig{ZB\_AM\_LS}\hfill\schemefig{ZB\_AM\_LSM}

\medskip\noindent
\schemefig{ZB\_AMISO\_LS}\hfill\schemefig{ZB\_AMISO\_LSM}

\end{document}
